% GENERATED FILE. Edit canonical lessons, then run npm run generate:books.
% canonical-source-hash: fnv1a64:694e32167e3e545d
% canonical-lessons: HI-C326-khandhar, HI-C326-khodna, HI-C326-murti, HI-C326-silalekh, HI-C326-avses

\chapter{Ruins of the Past}
\label{ch:hi-ruins-of-the-past}
\hlchaptermodality{326}

\begin{chapteropening}
\textbf{By the end of this chapter:} \emph{I can describe old ruins and what was found there.}

Complete the last lesson of chapter 326: I can describe old ruins and what was found there.
\end{chapteropening}

\section[\texorpdfstring{\hi{खंडहर} (khaṇḍhar)}{खंडहर (khaṇḍhar)}]{\hi{खंडहर} (khaṇḍhar) — ruins}

\label{lesson:HI-C326-khandhar}



\begin{quote}
Before the new one: say the Hindi for a sacrifice, then the Hindi for to entrust.
\end{quote}

\subsection*{You'll want to know: \hi{खंडहर}}
\textbf{\hi{खंडहर}} — \emph{khaṇḍhar} — \textquotedblleft{}ruins\textquotedblright{}.

\begin{quote}
\hi{खंडहर} \hi{के} \hi{लिए} \hi{सीधे} \hi{जाइए} \hi{और} \hi{मंदिर} \hi{के} \hi{बाद} \hi{दाएँ} \hi{मुड़िए।}
\end{quote}

\begin{quote}
\emph{khaṇḍhar ke liye sīdhe jāie aur mandir ke bād dāẽ muṛiye.}
\end{quote}

\begin{quote}
For the ruins, go straight on and turn right after the temple.
\end{quote}

Masculine.

\begin{cousinweb}
\textbf{\hi{खंडहर}} (\emph{khaṇḍhar}) is probably \emph{khaṇḍ}, broken, from Sanskrit \emph{khaṇḍa}, \textquotedblleft{}a piece, a fragment\textquotedblright{} (Monier-Williams), with \textbf{\hi{घर}}, a house, worn down: a house in pieces.
\end{cousinweb}

\begin{grammarlens}[title={What you've built}]
One more word for telling what happened.
\end{grammarlens}

\subsection*{Your turn}
\begin{itemize}
\raggedright
  \item \emph{Say it:} \emph{khaṇḍhar}
  \item \emph{Say it:} \emph{khaṇḍhar}, once more
  \item \emph{Say it:} \emph{khaṇḍhar ke liye sīdhe jāie aur mandir ke bād dāẽ muṛiye.} — For the ruins, go straight on and turn right after the temple.
  \item \emph{Recall:} say the Hindi for a sacrifice, then the Hindi for to entrust, then say \emph{khaṇḍhar} again
\end{itemize}

\begin{tcolorbox}[breakable,colback=teal!4,colframe=teal!35!black,title={Before you move on}]
What does \hi{खंडहर} mean? (\textquotedblleft{}Ruins\textquotedblright{}.) Say it once more.
\end{tcolorbox}
\section[\texorpdfstring{\hi{खोदना} (khodnā)}{खोदना (khodnā)}]{\hi{खोदना} (khodnā) — to dig}

\label{lesson:HI-C326-khodna}



\begin{quote}
Before the new one: say the Hindi for to entrust, then the Hindi for ruins.
\end{quote}

\subsection*{You'll want to know: \hi{खोदना}}
\textbf{\hi{खोदना}} — \emph{khodnā} — \textquotedblleft{}to dig\textquotedblright{}.

\begin{quote}
\hi{मज़दूरों} \hi{ने} \hi{खंडहर} \hi{के} \hi{पास} \hi{मिट्टी} \hi{खोदी।}
\end{quote}

\begin{quote}
\emph{mazdūroṁ ne khaṇḍhar ke pās miṭṭī khodī.}
\end{quote}

\begin{quote}
The labourers dug up the earth near the ruins.
\end{quote}

Past, with \hi{ने} (he or she …): \hi{उसने} \hi{कुआँ} \hi{खोदा}, \emph{usne kuā̃ khodā}.

\begin{cousinweb}
\textbf{\hi{खोदना}} (\emph{khodnā}) is inherited from \emph{*khōdd}, a form CDIAL reconstructs for digging (CDIAL 3934). What is dug agrees: \textbf{\hi{मिट्टी} \hi{खोदी}}, \textbf{\hi{कुआँ} \hi{खोदा}}.
\end{cousinweb}

\begin{grammarlens}[title={What you've built}]
One more word for telling what happened.
\end{grammarlens}

\subsection*{Your turn}
\begin{itemize}
\raggedright
  \item \emph{Say it:} \emph{khodnā}
  \item \emph{Say it:} \emph{khodnā}, once more
  \item \emph{Say it:} \emph{mazdūroṁ ne khaṇḍhar ke pās miṭṭī khodī.} — The labourers dug up the earth near the ruins.
  \item \emph{Recall:} say the Hindi for to entrust, then the Hindi for ruins, then say \emph{khodnā} again
\end{itemize}

\begin{tcolorbox}[breakable,colback=teal!4,colframe=teal!35!black,title={Before you move on}]
What does \hi{खोदना} mean? (\textquotedblleft{}To dig\textquotedblright{}.) Say it once more.
\end{tcolorbox}
\section[\texorpdfstring{\hi{मूर्ति} (mūrti)}{मूर्ति (mūrti)}]{\hi{मूर्ति} (mūrti) — a statue; an image}

\label{lesson:HI-C326-murti}



\begin{quote}
Before the new one: say the Hindi for ruins, then the Hindi for to dig.
\end{quote}

\subsection*{You'll want to know: \hi{मूर्ति}}
\textbf{\hi{मूर्ति}} — \emph{mūrti} — \textquotedblleft{}a statue; an image\textquotedblright{}.

\begin{quote}
\hi{मज़दूरों} \hi{को} \hi{मिट्टी} \hi{में} \hi{एक} \hi{पत्थर} \hi{की} \hi{मूर्ति} \hi{मिली।}
\end{quote}

\begin{quote}
\emph{mazdūroṁ ko miṭṭī meṁ ek patthar kī mūrti milī.}
\end{quote}

\begin{quote}
The labourers found a stone statue in the earth.
\end{quote}

Feminine.

\begin{cousinweb}
\textbf{\hi{मूर्ति}} (\emph{mūrti}) is Sanskrit \emph{mūrti}, \textquotedblleft{}a solid body, a form\textquotedblright{} (Monier-Williams): a figure of stone, metal or clay, and in a temple the image of a god.
\end{cousinweb}

\begin{grammarlens}[title={What you've built}]
One more word for telling what happened.
\end{grammarlens}

\subsection*{Your turn}
\begin{itemize}
\raggedright
  \item \emph{Say it:} \emph{mūrti}
  \item \emph{Say it:} \emph{mūrti}, once more
  \item \emph{Say it:} \emph{mazdūroṁ ko miṭṭī meṁ ek patthar kī mūrti milī.} — The labourers found a stone statue in the earth.
  \item \emph{Recall:} say the Hindi for ruins, then the Hindi for to dig, then say \emph{mūrti} again
\end{itemize}

\begin{tcolorbox}[breakable,colback=teal!4,colframe=teal!35!black,title={Before you move on}]
What does \hi{मूर्ति} mean? (\textquotedblleft{}A statue; an image\textquotedblright{}.) Say it once more.
\end{tcolorbox}
\section[\texorpdfstring{\hi{शिलालेख} (śilālekh)}{शिलालेख (śilālekh)}]{\hi{शिलालेख} (śilālekh) — an inscription (on stone)}

\label{lesson:HI-C326-silalekh}



\begin{quote}
Before the new one: say the Hindi for to dig, then the Hindi for a statue.
\end{quote}

\subsection*{You'll want to know: \hi{शिलालेख}}
\textbf{\hi{शिलालेख}} — \emph{śilālekh} — \textquotedblleft{}an inscription (on stone)\textquotedblright{}.

\begin{quote}
\hi{उस} \hi{शिलालेख} \hi{पर} \hi{एक} \hi{राजा} \hi{का} \hi{नाम} \hi{है।}
\end{quote}

\begin{quote}
\emph{us śilālekh par ek rājā kā nām hai.}
\end{quote}

\begin{quote}
That inscription bears a king's name.
\end{quote}

Masculine.

\begin{cousinweb}
\textbf{\hi{शिलालेख}} (\emph{śilālekh}) joins Sanskrit \emph{śilā}, \textquotedblleft{}a stone\textquotedblright{}, and \emph{lekha}, \textquotedblleft{}a writing\textquotedblright{} (Monier-Williams): words cut into rock, one of the oldest sources of history.
\end{cousinweb}

\begin{grammarlens}[title={What you've built}]
One more word for telling what happened.
\end{grammarlens}

\subsection*{Your turn}
\begin{itemize}
\raggedright
  \item \emph{Say it:} \emph{śilālekh}
  \item \emph{Say it:} \emph{śilālekh}, once more
  \item \emph{Say it:} \emph{us śilālekh par ek rājā kā nām hai.} — That inscription bears a king's name.
  \item \emph{Recall:} say the Hindi for to dig, then the Hindi for a statue, then say \emph{śilālekh} again
\end{itemize}

\begin{tcolorbox}[breakable,colback=teal!4,colframe=teal!35!black,title={Before you move on}]
What does \hi{शिलालेख} mean? (\textquotedblleft{}An inscription (on stone)\textquotedblright{}.) Say it once more.
\end{tcolorbox}
\section[\texorpdfstring{\hi{अवशेष} (avśeṣ)}{अवशेष (avśeṣ)}]{\hi{अवशेष} (avśeṣ) — remains}

\label{lesson:HI-C326-avses}



\begin{quote}
Before the new one: say the Hindi for a statue, then the Hindi for an inscription (on stone).
\end{quote}

\subsection*{You'll want to know: \hi{अवशेष}}
\textbf{\hi{अवशेष}} — \emph{avśeṣ} — \textquotedblleft{}remains\textquotedblright{}.

\begin{quote}
\hi{खंडहर} \hi{में} \hi{एक} \hi{पुराने} \hi{मंदिर} \hi{के} \hi{अवशेष} \hi{भी} \hi{हैं।}
\end{quote}

\begin{quote}
\emph{khaṇḍhar meṁ ek purāne mandir ke avśeṣ bhī haiṁ.}
\end{quote}

\begin{quote}
Among the ruins are the remains of an old temple too.
\end{quote}

Masculine.

\begin{cousinweb}
\textbf{\hi{अवशेष}} (\emph{avśeṣ}) is Sanskrit \emph{avaśeṣa}, \textquotedblleft{}a remainder, what is left\textquotedblright{} (Monier-Williams). It is mostly used in the plural sense: what is left of a building or a body.
\end{cousinweb}

\begin{grammarlens}[title={What you've built}]
One more word for telling what happened.
\end{grammarlens}

\subsection*{Your turn}
\begin{itemize}
\raggedright
  \item \emph{Say it:} \emph{avśeṣ}
  \item \emph{Say it:} \emph{avśeṣ}, once more
  \item \emph{Say it:} \emph{khaṇḍhar meṁ ek purāne mandir ke avśeṣ bhī haiṁ.} — Among the ruins are the remains of an old temple too.
  \item \emph{Recall:} say the Hindi for a statue, then the Hindi for an inscription (on stone), then say \emph{avśeṣ} again
\end{itemize}

\begin{tcolorbox}[breakable,colback=teal!4,colframe=teal!35!black,title={Before you move on}]
What does \hi{अवशेष} mean? (\textquotedblleft{}Remains\textquotedblright{}.) Say it once more.
\end{tcolorbox}
